% ============================================================================
%  Inverse Kaluza-Klein Scale Symmetry Framework
%  Document II: Initial Draft
%  Large independent import sections in TeX.
%  Independent chains of equations and code.
%  Compile with pdflatex (twice for refs).
% ============================================================================
\documentclass[11pt, a4paper, oneside]{article}

% --- Geometry ----------------------------------------------------------------
\usepackage[
  top=1.0in, bottom=1.0in, left=1.15in, right=1.15in,
  headheight=14pt
]{geometry}

% --- Fonts / encoding --------------------------------------------------------
\usepackage[T1]{fontenc}
\usepackage{lmodern}
\usepackage{microtype}

% --- Math --------------------------------------------------------------------
\usepackage{amsmath, amssymb, amsthm, mathtools}
\usepackage{bm}
\usepackage{physics}

% --- Code listings -----------------------------------------------------------
\usepackage{listings}
\usepackage{xcolor}
\definecolor{codegray}{rgb}{0.5,0.5,0.5}
\definecolor{codeblue}{rgb}{0.13,0.13,0.9}
\definecolor{codegreen}{rgb}{0.0,0.5,0.0}
\definecolor{codebg}{rgb}{0.97,0.97,0.97}
\lstset{
  backgroundcolor=\color{codebg},
  basicstyle=\ttfamily\small,
  keywordstyle=\color{codeblue}\bfseries,
  commentstyle=\color{codegreen}\itshape,
  stringstyle=\color{orange},
  numberstyle=\tiny\color{codegray},
  numbers=left,
  numbersep=8pt,
  breaklines=true,
  frame=single,
  rulecolor=\color{codegray},
  captionpos=b,
  showstringspaces=false,
  tabsize=2
}

% --- Tables ------------------------------------------------------------------
\usepackage{booktabs}
\usepackage{array}
\usepackage{tabularx}

% --- Graphics (for diagrams) -------------------------------------------------
\usepackage{tikz}
\usetikzlibrary{arrows.meta, positioning, shapes.geometric, fit, calc}

% --- Lists -------------------------------------------------------------------
\usepackage{enumitem}
\setlist{nosep, leftmargin=1.5em}

% --- Headers / footers -------------------------------------------------------
\usepackage{fancyhdr}
\pagestyle{fancy}
\fancyhf{}
\fancyhead[L]{\small\textit{Inverse KK Scale Symmetry}}
\fancyhead[R]{\small\textit{Doc II --- Initial Draft}}
\fancyfoot[C]{\thepage}
\renewcommand{\headrulewidth}{0.4pt}

% --- Section formatting ------------------------------------------------------
\usepackage{titlesec}
\titleformat{\section}
  {\Large\bfseries}
  {\S\thesection}{1em}{}
\titleformat{\subsection}
  {\large\bfseries}
  {\thesubsection}{1em}{}

% --- Hyperlinks --------------------------------------------------------------
\usepackage[
  colorlinks=true, linkcolor=black, citecolor=black, urlcolor=blue
]{hyperref}

% --- Theorem-like environments -----------------------------------------------
\theoremstyle{definition}
\newtheorem{axiom}{Axiom}[section]
\newtheorem{definition}{Definition}[section]
\theoremstyle{plain}
\newtheorem{theorem}{Theorem}[section]
\newtheorem{proposition}{Proposition}[section]
\theoremstyle{remark}
\newtheorem{remark}{Remark}[section]

% --- Custom operators --------------------------------------------------------
\DeclareMathOperator{\Pack}{\Pi}
\newcommand{\Id}{\mathcal{I}}
\newcommand{\Res}{\mathcal{R}}
\newcommand{\Proj}{P}

% ============================================================================
\begin{document}

\begin{titlepage}
  \centering
  \vspace*{2cm}
  {\Huge\bfseries Inverse Kaluza-Klein Scale Symmetry\par}
  \vspace{0.8cm}
  {\LARGE Document II --- Initial Draft\par}
  \vspace{0.5cm}
  {\Large Independent import sections, equation chains, and code\par}
  \vspace{1.5cm}
  {\normalsize
	This document is structured as a collection of independent sections,
	each self-contained. Each section may be read, compiled, or imported
	separately. Order is thematic, not logically required.
  \par}
  \vspace{1.5cm}
  {\large VSEPR-SIM Theoretical Division\par}
  {\large Development Day 64\par}
\end{titlepage}

\tableofcontents
\newpage

% ============================================================================
%  IMPORT SECTION A
%  Title: Framework Identity and Naming
%  Status: Standalone
% ============================================================================
\section{Framework Identity and Naming}
\label{sec:naming}

% --- Begin standalone import block ------------------------------------------
\noindent\rule{\linewidth}{0.8pt}
\textbf{Import block A} \hfill \textit{standalone, no dependencies}
\noindent\rule{\linewidth}{0.4pt}

\subsection*{The Direction of Inference}

Standard Kaluza-Klein (KK) logic:
\begin{equation}
  \text{extra dimension}
  \;\xrightarrow{\text{KK}}\;
  \text{force behavior}
  \tag{A.1}
\end{equation}

Inverse KK logic (this framework):
\begin{equation}
  \boxed{
	\text{observed force-scale behavior}
	\;\xrightarrow{\text{IKK}}\;
	\text{implied dimensional permission}
  }
  \tag{A.2}
\end{equation}

The framework does not assume extra dimensions and then derive consequences.
It starts from measured force-scale residuals and asks what dimensional
permission structure would \emph{produce} those residuals as projections.

\subsection*{Why ``Inverse''}

The word ``inverse'' is not metaphorical. The mathematical arrow is
literally reversed:
\begin{center}
  \begin{tikzpicture}[node distance=1.5cm, auto,
	  box/.style={draw, rectangle, rounded corners, minimum width=3.5cm,
				  minimum height=0.8cm, text centered}]
	\node[box] (D) {Extra dimension $D$};
	\node[box, right=2cm of D] (F) {Force behavior $F$};
	\draw[->, thick, blue] (D) -- node[above]{\small KK} (F);
	\draw[->, thick, red, dashed] (F) -- node[below]{\small IKK (this work)} (D);
  \end{tikzpicture}
\end{center}

\noindent\rule{\linewidth}{0.4pt}
% --- End standalone import block --------------------------------------------

% ============================================================================
%  IMPORT SECTION B
%  Title: Scale Ladder --- Axiom Chain
%  Status: Standalone
% ============================================================================
\section{Scale Ladder: Axiom Chain}
\label{sec:axioms}

\noindent\rule{\linewidth}{0.8pt}
\textbf{Import block B} \hfill \textit{standalone, no dependencies}
\noindent\rule{\linewidth}{0.4pt}

\begin{axiom}[Existence gate]
  Every physical entity satisfies a binary existence condition:
  \begin{equation}
	E \in \{0, 1\}
	\qquad S_0 = (\text{existence},\; \text{---},\; \lambda_\text{Planck})
	\tag{B.1}
  \end{equation}
\end{axiom}

\begin{axiom}[Scale entry]
  Each scale level adds one degree of freedom, one mediator, and one
  observational limit:
  \begin{equation}
	S_n = (D_n,\; M_n,\; L_n)
	\tag{B.2}
  \end{equation}
\end{axiom}

\begin{axiom}[Packing occurs]
  Components at scale $n$ self-assemble into identities at scale $n+1$:
  \begin{equation}
	\mathbf{C}^{(n)} \xrightarrow{\;\Pi_n\;} \mathbf{I}^{(n+1)}
	\tag{B.3}
  \end{equation}
  Packing is irreversible without energy cost $\geq E_\text{bind}$.
\end{axiom}

\begin{axiom}[Projection]
  Observation at scale $S$ produces only:
  \begin{equation}
	\mathcal{O}_S = P_S(\mathcal{I})
	\tag{B.4}
  \end{equation}
  The full identity $\mathcal{I}$ is not directly accessible.
\end{axiom}

\begin{axiom}[Residual existence]
  The residual is the difference between the full identity content at
  scale $S$ and what observation recovers:
  \begin{equation}
	\mathcal{R}_S = \mathcal{U}_S(\mathcal{I}) - \mathcal{E}_S(\mathcal{O}_S)
	\tag{B.5}
  \end{equation}
  $\mathcal{R}_S = 0$ only if projection is perfectly injective at scale $S$.
\end{axiom}

\noindent\rule{\linewidth}{0.4pt}

% ============================================================================
%  IMPORT SECTION C
%  Title: Matrix Packing Chain
%  Status: Standalone
% ============================================================================
\section{Matrix Packing Chain}
\label{sec:packing}

\noindent\rule{\linewidth}{0.8pt}
\textbf{Import block C} \hfill \textit{standalone, no dependencies}
\noindent\rule{\linewidth}{0.4pt}

\subsection*{C.1 Quark to Particle}

Quark component matrix for particle $p$:
\begin{equation}
  \mathbf{C}_q =
  \begin{bmatrix}
	q_1 & c_1 & \eta_1 \\
	q_2 & c_2 & \eta_2 \\
	q_3 & c_3 & \eta_3
  \end{bmatrix}
  \in \mathbb{R}^{3\times 3}
  \tag{C.1}
\end{equation}
where columns are charge, color, and helicity channels.

Packing to particle identity row-vector:
\begin{equation}
  \mathbf{I}_p = \begin{bmatrix} Q_p & M_p & S_p \end{bmatrix}
  \in \mathbb{R}^{1\times 3}
  \tag{C.2}
\end{equation}

Formal map:
\begin{equation}
  \mathbf{I}_p = \Pi_q(\mathbf{C}_q)
  \qquad 3\times 3 \to 1\times 3
  \tag{C.3}
\end{equation}

\subsection*{C.2 Recursive Packing Ladder}

\begin{equation}
  \underbrace{\text{quarks}}_{S_2}
  \xrightarrow{\Pi_2}
  \underbrace{\text{particles}}_{S_3}
  \xrightarrow{\Pi_3}
  \underbrace{\text{atoms}}_{S_3}
  \xrightarrow{\Pi_3}
  \underbrace{\text{molecules}}_{S_3}
  \xrightarrow{\Pi_4}
  \underbrace{\text{materials}}_{S_4}
  \xrightarrow{\Pi_5}
  \underbrace{\text{grav.\ systems}}_{S_5}
  \xrightarrow{\Pi_D}
  \underbrace{\text{dark residuals}}_{S_D}
  \tag{C.4}
\end{equation}

\subsection*{C.3 Entropy as Packing Degeneracy}

The entropy of a packed identity at scale $n$:
\begin{equation}
  \mathcal{S}_n = k_B \ln \Omega(\mathbf{I}^{(n)})
  \tag{C.5}
\end{equation}
where $\Omega(\mathbf{I}^{(n)})$ counts the number of component
configurations $\mathbf{C}^{(n-1)}$ that map to the same
$\mathbf{I}^{(n)}$ under $\Pi_{n-1}$.

\subsection*{C.4 Complexity}

Stable packed identity plus structured residual:
\begin{equation}
  \mathcal{C}_n = f\!\left(\mathbf{I}^{(n+1)},\; \mathcal{R}_n\right)
  \tag{C.6}
\end{equation}

High complexity requires both a non-trivial packed identity and a
non-trivial residual. Pure structure with zero residual is simple.
Pure noise with no packed identity is also simple. Complexity is the
overlap.

\noindent\rule{\linewidth}{0.4pt}

% ============================================================================
%  IMPORT SECTION D
%  Title: Cancellation and Hidden Intensity
%  Status: Standalone
% ============================================================================
\section{Cancellation and Hidden Intensity}
\label{sec:cancellation}

\noindent\rule{\linewidth}{0.8pt}
\textbf{Import block D} \hfill \textit{standalone, no dependencies}
\noindent\rule{\linewidth}{0.4pt}

\subsection*{D.1 The Rule}

\begin{equation}
  \boxed{
	\text{net-zero packed value}
	\;\neq\;
	\text{no internal structure}
  }
  \tag{D.1}
\end{equation}

\subsection*{D.2 Observable and Hidden Channels}

For a component vector $\mathbf{c}_j \in \mathbb{R}^m$:
\begin{align}
  I_j &= \bm{\alpha}^\top \mathbf{c}_j
  \tag{D.2} \\
  H_j &= \mathbf{c}_j^\top \mathbf{W}\, \mathbf{c}_j
  \tag{D.3}
\end{align}
where $\bm{\alpha}$ is the projection weight vector and $\mathbf{W}$ is a
positive-semi-definite weight matrix for the hidden channel.

Hidden-active condition:
\begin{equation}
  I_j = 0, \quad H_j > 0
  \tag{D.4}
\end{equation}

\subsection*{D.3 Neutron-Like Example}

Charge projection ($\bm{\alpha} = [1, 1, 1]^\top$):
\begin{equation}
  I_\text{charge}
  = +\tfrac{2}{3} + (-\tfrac{1}{3}) + (-\tfrac{1}{3})
  = 0
  \tag{D.5}
\end{equation}

Diagonal hidden intensity ($\mathbf{W} = \mathbf{I}_{3\times 3}$):
\begin{equation}
  H_\text{charge}
  = \left(\tfrac{2}{3}\right)^2
  + \left(\tfrac{1}{3}\right)^2
  + \left(\tfrac{1}{3}\right)^2
  = \tfrac{4}{9} + \tfrac{1}{9} + \tfrac{1}{9}
  = \tfrac{6}{9}
  = \tfrac{2}{3}
  \neq 0
  \tag{D.6}
\end{equation}

\noindent\rule{\linewidth}{0.4pt}

% ============================================================================
%  IMPORT SECTION E
%  Title: Dark Matter Kernel Equations
%  Status: Standalone
% ============================================================================
\section{Dark Matter Kernel Equations}
\label{sec:dark}

\noindent\rule{\linewidth}{0.8pt}
\textbf{Import block E} \hfill \textit{standalone, no dependencies}
\noindent\rule{\linewidth}{0.4pt}

\subsection*{E.1 Projection Gap Definition}

\begin{align}
  \rho_\text{dark}
  &= \rho_G - \rho_\text{EM}
  \tag{E.1} \\
  &= P_G(\rho) - P_\text{EM}(\rho)
  \tag{E.2}
\end{align}

\subsection*{E.2 Kernel Form}

\begin{equation}
  \boxed{
	\rho_\text{dark}(x,y,z;\,t)
	= \int \rho(x,y,z,w;\,t)
	  \Bigl[K_G(w) - K_\text{EM}(w)\Bigr]\, dw
  }
  \tag{E.3}
\end{equation}

\subsection*{E.3 Kernel Properties (assumed)}

Gravitational kernel: couples to all scales
\begin{equation}
  \int K_G(w)\, dw = 1
  \tag{E.4}
\end{equation}

EM kernel: truncates above $L_\text{EM}$
\begin{equation}
  K_\text{EM}(w) \approx 0 \quad \text{for } w > L_\text{EM}
  \tag{E.5}
\end{equation}

The integrand $[K_G(w) - K_\text{EM}(w)]$ is therefore non-zero precisely
in the dark regime $w > L_\text{EM}$.

\subsection*{E.4 Dark fraction}

\begin{equation}
  f_\text{dark} = \frac{\int_{L_\text{EM}}^\infty
	\rho(x,y,z,w;\,t)\, K_G(w)\, dw}
	{\int_0^\infty \rho(x,y,z,w;\,t)\, K_G(w)\, dw}
  \tag{E.6}
\end{equation}

Observed dark-matter fraction from cosmological measurements:
$f_\text{dark} \approx 0.27$.

\noindent\rule{\linewidth}{0.4pt}

% ============================================================================
%  IMPORT SECTION F
%  Title: Effective Proper Time Chain
%  Status: Standalone
% ============================================================================
\section{Effective Proper Time Chain}
\label{sec:propertime}

\noindent\rule{\linewidth}{0.8pt}
\textbf{Import block F} \hfill \textit{standalone; references import block B}
\noindent\rule{\linewidth}{0.4pt}

\subsection*{F.1 Classical baseline}

\begin{equation}
  d\tau_\text{4D}
  = dt\,\sqrt{
	  1 - \frac{v^2}{c^2} + \frac{2\Phi_G}{c^2}
	}
  \tag{F.1}
\end{equation}

Integrated:
\begin{equation}
  \tau_\text{4D}
  = \int_0^T dt\,\sqrt{
	  1 - \frac{v^2(t)}{c^2} + \frac{2\Phi_G(t)}{c^2}
	}
  \tag{F.2}
\end{equation}

\subsection*{F.2 Quantum proper time (ion clock result)}

\begin{equation}
  \hat{\tau} \equiv \tau(\hat{x}, \hat{p})
  \tag{F.3}
\end{equation}

Proper time becomes quantum-linked. The observable has spread:
\begin{equation}
  R_\tau = \hat{\tau} - \langle\hat{\tau}\rangle
  \tag{F.4}
\end{equation}

\subsection*{F.3 Dark-scale extension}

\begin{equation}
  \boxed{
	d\tau_\text{eff}
	= dt\,\sqrt{
		1 - \frac{v^2}{c^2} + \frac{2\Phi_G}{c^2} + \chi_w(w)
	  }
  }
  \tag{F.5}
\end{equation}

Additive decomposition:
\begin{equation}
  d\tau_\text{eff} = d\tau_\text{4D} + d\tau_w
  \tag{F.6}
\end{equation}

Limit check:
\begin{equation}
  \chi_w(w) \to 0 \implies d\tau_\text{eff} \to d\tau_\text{4D}
  \tag{F.7}
\end{equation}

\subsection*{F.4 Proper-time identity bundle}

\begin{equation}
  \mathcal{I}_\tau
  = \bigl(\hat{x},\; \hat{p},\; \hat{\tau},\; H_c,\; H_m,\; w\bigr)
  \tag{F.8}
\end{equation}

Clock-motion relation particle:
\begin{equation}
  \mathcal{R}_{cm} = [0,\, 0,\, 0,\, 0 \mid \mathbf{T}_{cm}]
  \tag{F.9}
\end{equation}

\noindent\rule{\linewidth}{0.4pt}

% ============================================================================
%  IMPORT SECTION G
%  Title: Entanglement Chain
%  Status: Standalone
% ============================================================================
\section{Entanglement: Relation-Particle Chain}
\label{sec:entanglement}

\noindent\rule{\linewidth}{0.8pt}
\textbf{Import block G} \hfill \textit{standalone, no dependencies}
\noindent\rule{\linewidth}{0.4pt}

\subsection*{G.1 Relation object}

\begin{equation}
  \mathcal{R}_{AB}
  = \bigl[\underbrace{0,\; 0,\; 0,\; 0}_{\Delta x,\Delta y,\Delta z,\Delta t}
	\;\big|\;
	\mathbf{T}_{AB}\bigr]
  \tag{G.1}
\end{equation}

\begin{equation}
  \mathbf{T}_{AB} \in \{-1,\, 0,\, +1\}^m
  \tag{G.2}
\end{equation}

\subsection*{G.2 Pair decomposition}

\begin{equation}
  \mathcal{I}_{AB}
  \;\to\;
  \mathcal{I}_A + \mathcal{I}_B + \mathcal{R}_{AB}
  \tag{G.3}
\end{equation}

\subsection*{G.3 Measurement}

\begin{align}
  O_A &= P_a(\mathcal{I}_A,\, \mathcal{R}_{AB}) \tag{G.4} \\
  O_B &= P_b(\mathcal{I}_B,\, \mathcal{R}_{AB}) \tag{G.5} \\
  E(a,b) &= \langle O_A O_B \rangle \tag{G.6}
\end{align}

\subsection*{G.4 No-signal interpretation}

$\mathcal{R}_{AB}$ is a fixed structure in identity-space. It does not
propagate. Measurement at $A$ reads the relation; measurement at $B$
reads the same relation. No signal travels.

\begin{equation}
  \text{real} \neq \text{spatially displaced}
  \quad,\qquad
  \text{deterministic} \neq \text{time-evolving signal}
  \tag{G.7}
\end{equation}

\noindent\rule{\linewidth}{0.4pt}

% ============================================================================
%  IMPORT SECTION H
%  Title: Code --- Python numerical prototype
%  Status: Standalone
% ============================================================================
\section{Code: Python Numerical Prototype}
\label{sec:code-python}

\noindent\rule{\linewidth}{0.8pt}
\textbf{Import block H} \hfill \textit{standalone, Python 3.10+}
\noindent\rule{\linewidth}{0.4pt}

The following code implements the core algebraic objects in a minimal
self-contained Python prototype. No physics library is required.

\begin{lstlisting}[language=Python, caption={Identity bundle and projection}]
import numpy as np
from dataclasses import dataclass, field
from typing import Optional

# --- Scale ladder entry ------------------------------------------------------

@dataclass
class ScaleEntry:
	n:      int
	label:  str
	dof:    str          # degree of freedom name
	mediator: str        # force carrier / mediator name
	L_m:    float        # observational limit (metres)

SCALE_LADDER = [
	ScaleEntry(0, "S0", "existence",   "---",      1.616e-35),
	ScaleEntry(1, "S1", "vector",      "---",      1e-19),
	ScaleEntry(2, "S2", "angular",     "gluon",    1e-15),
	ScaleEntry(3, "S3", "spatial",     "photon",   1e-10),
	ScaleEntry(4, "S4", "time",        "W/Z",      1e-18),
	ScaleEntry(5, "S5", "curvature",   "graviton", 1e26),
	ScaleEntry(6, "SD", "dark",        "X_D",      float('inf')),
]

# --- Identity component ------------------------------------------------------

@dataclass
class Component:
	"""One component of an identity bundle at a given scale."""
	scale_n:  int
	value:    np.ndarray   # shape (m,)  -- multi-channel
	weight:   float = 1.0  # alpha_k for anchor computation

# --- Identity bundle ---------------------------------------------------------

@dataclass
class IdentityBundle:
	"""Full identity state I_i(t)."""
	label:      str
	components: list[Component] = field(default_factory=list)

	def anchor(self) -> np.ndarray:
		"""
		Identity anchor A_i = (sum alpha_k * c_k) / (sum alpha_k)
		Returns weighted mean of component values.
		"""
		total_w = sum(c.weight for c in self.components)
		if total_w == 0.0:
			raise ValueError("Zero total weight in anchor computation.")
		return sum(c.weight * c.value for c in self.components) / total_w

	def fuzz_radius(self, proj_indices=(0, 1, 2)) -> float:
		"""
		Fuzz radius at 3D spatial projection.
		proj_indices: which component channels correspond to x,y,z.
		"""
		r_obs = self.anchor()[list(proj_indices)]
		total_w = sum(c.weight for c in self.components)
		sq_sum = sum(
			c.weight * np.sum((c.value[list(proj_indices)] - r_obs)**2)
			for c in self.components
		)
		return np.sqrt(sq_sum / total_w)

# --- Projection operator -----------------------------------------------------

def project_EM(bundle: IdentityBundle, L_EM: float) -> np.ndarray:
	"""
	P_EM: retain only components at scales with L_m <= L_EM.
	Returns the projected anchor value.
	"""
	em_components = [
		c for c in bundle.components
		if SCALE_LADDER[c.scale_n].L_m <= L_EM
	]
	if not em_components:
		return np.zeros_like(bundle.components[0].value)
	total_w = sum(c.weight for c in em_components)
	return sum(c.weight * c.value for c in em_components) / total_w

# --- Dark density residual ---------------------------------------------------

def dark_density_residual(
	rho_full: np.ndarray,   # shape (Nw,)  density as function of w
	w_grid:   np.ndarray,   # shape (Nw,)  w coordinate values
	K_G:      np.ndarray,   # shape (Nw,)  gravitational kernel
	K_EM:     np.ndarray,   # shape (Nw,)  EM kernel
) -> float:
	"""
	rho_dark = integral rho(w) [K_G(w) - K_EM(w)] dw
	Numerical trapezoidal integration.
	"""
	integrand = rho_full * (K_G - K_EM)
	return float(np.trapz(integrand, w_grid))

# --- Cancellation and hidden intensity ---------------------------------------

def observable_channel(c_j: np.ndarray, alpha: np.ndarray) -> float:
	"""I_j = alpha^T c_j"""
	return float(alpha @ c_j)

def hidden_intensity(c_j: np.ndarray, W: np.ndarray) -> float:
	"""H_j = c_j^T W c_j"""
	return float(c_j @ W @ c_j)

def is_hidden_active(
	c_j: np.ndarray,
	alpha: np.ndarray,
	W: np.ndarray,
	tol: float = 1e-12
) -> bool:
	"""Returns True when I_j == 0 but H_j > 0."""
	I = abs(observable_channel(c_j, alpha))
	H = hidden_intensity(c_j, W)
	return (I < tol) and (H > tol)

# --- Effective proper time ---------------------------------------------------

def dtau_eff(
	dt:     float,
	v:      float,
	c:      float = 2.998e8,
	Phi_G:  float = 0.0,
	chi_w:  float = 0.0,
) -> float:
	"""
	d tau_eff = dt * sqrt(1 - v^2/c^2 + 2*Phi_G/c^2 + chi_w)
	chi_w: dark-scale metric correction (default 0 -> standard GR).
	"""
	arg = 1.0 - (v/c)**2 + 2.0*Phi_G/c**2 + chi_w
	if arg < 0.0:
		raise ValueError(f"Negative radicand in dtau_eff: {arg:.4e}")
	return dt * np.sqrt(arg)

# --- Relation particle -------------------------------------------------------

@dataclass
class RelationParticle:
	"""
	R_AB = [0, 0, 0, 0 | T_AB]
	Spacetime displacement is zero. Relation vector T carries the content.
	"""
	T: np.ndarray  # shape (m,), values in {-1, 0, +1}

	@property
	def spacetime_displacement(self) -> np.ndarray:
		return np.zeros(4)

	def measure_A(self, I_A: IdentityBundle, projection_A) -> np.ndarray:
		return projection_A(I_A, self)

	def measure_B(self, I_B: IdentityBundle, projection_B) -> np.ndarray:
		return projection_B(I_B, self)
\end{lstlisting}

\noindent\rule{\linewidth}{0.4pt}

% ============================================================================
%  IMPORT SECTION I
%  Title: Code --- VSIM kernel hook sketch (C++)
%  Status: Standalone
% ============================================================================
\section{Code: VSIM Kernel Hook Sketch (C++)}
\label{sec:code-cpp}

\noindent\rule{\linewidth}{0.8pt}
\textbf{Import block I} \hfill \textit{standalone, C++23}
\noindent\rule{\linewidth}{0.4pt}

The following sketch shows how the theoretical objects above would be
wired into the VSIM kernel event system (see WO-64A / WO-64B).

\begin{lstlisting}[language=C++, caption={Identity anchor and hidden intensity
  in VSIM kernel style}]
// ikk_kernel_hook.hpp  --  WO-64 theoretical integration sketch
// Compile as part of include/core/
#pragma once

#include <array>
#include <cmath>
#include <numeric>
#include <span>
#include <stdexcept>
#include <vector>

namespace ikk {

// ---- Scale entry -----------------------------------------------------------

struct ScaleEntry {
	int         n;
	const char* label;
	const char* dof;
	const char* mediator;
	double      L_m;     // observational limit (m)
};

inline constexpr std::array<ScaleEntry, 7> kScaleLadder = {{
	{0, "S0", "existence",  "---",     1.616e-35},
	{1, "S1", "vector",     "---",     1.0e-19  },
	{2, "S2", "angular",    "gluon",   1.0e-15  },
	{3, "S3", "spatial",    "photon",  1.0e-10  },
	{4, "S4", "time",       "W/Z",     1.0e-18  },
	{5, "S5", "curvature",  "graviton",1.0e+26  },
	{6, "SD", "dark",       "X_D",     1.0e+99  },
}};

// ---- Component -------------------------------------------------------------

struct Component {
	int                 scale_n;
	std::vector<double> value;    // multi-channel value
	double              weight = 1.0;
};

// ---- Identity anchor -------------------------------------------------------

std::vector<double> identity_anchor(
	std::span<const Component> components)
{
	if (components.empty())
		throw std::runtime_error("ikk: empty component set");

	const std::size_t m = components[0].value.size();
	std::vector<double> acc(m, 0.0);
	double total_w = 0.0;

	for (const auto& c : components) {
		total_w += c.weight;
		for (std::size_t k = 0; k < m; ++k)
			acc[k] += c.weight * c.value[k];
	}
	for (auto& v : acc) v /= total_w;
	return acc;
}

// ---- Fuzz radius -----------------------------------------------------------

double fuzz_radius(
	std::span<const Component> components,
	std::span<const std::size_t> proj_idx)   // which channels = 3D
{
	const auto anchor = identity_anchor(components);
	double sq_sum = 0.0, total_w = 0.0;

	for (const auto& c : components) {
		total_w += c.weight;
		double d2 = 0.0;
		for (auto idx : proj_idx) {
			double diff = c.value[idx] - anchor[idx];
			d2 += diff * diff;
		}
		sq_sum += c.weight * d2;
	}
	return std::sqrt(sq_sum / total_w);
}

// ---- Observable channel and hidden intensity --------------------------------

double observable_channel(
	std::span<const double> c_j,
	std::span<const double> alpha)
{
	return std::inner_product(c_j.begin(), c_j.end(),
							   alpha.begin(), 0.0);
}

// W is m x m, row-major
double hidden_intensity(
	std::span<const double> c_j,
	std::span<const double> W_flat,
	std::size_t m)
{
	double result = 0.0;
	for (std::size_t i = 0; i < m; ++i)
		for (std::size_t j = 0; j < m; ++j)
			result += c_j[i] * W_flat[i * m + j] * c_j[j];
	return result;
}

bool is_hidden_active(
	std::span<const double> c_j,
	std::span<const double> alpha,
	std::span<const double> W_flat,
	std::size_t m,
	double tol = 1.0e-12)
{
	const double I = std::abs(observable_channel(c_j, alpha));
	const double H = hidden_intensity(c_j, W_flat, m);
	return (I < tol) && (H > tol);
}

// ---- Effective proper time -------------------------------------------------

double dtau_eff(
	double dt,
	double v,
	double Phi_G  = 0.0,
	double chi_w  = 0.0,
	double c      = 2.99792458e8)
{
	const double arg = 1.0 - (v * v) / (c * c)
					 + 2.0 * Phi_G / (c * c)
					 + chi_w;
	if (arg < 0.0)
		throw std::runtime_error("ikk::dtau_eff: negative radicand");
	return dt * std::sqrt(arg);
}

// ---- Dark density residual -------------------------------------------------

double dark_density_residual(
	std::span<const double> rho_w,    // density as function of w
	std::span<const double> w_grid,   // w coordinate values
	std::span<const double> K_G,      // gravitational kernel
	std::span<const double> K_EM)     // EM kernel
{
	const std::size_t N = rho_w.size();
	std::vector<double> integrand(N);
	for (std::size_t i = 0; i < N; ++i)
		integrand[i] = rho_w[i] * (K_G[i] - K_EM[i]);

	// Trapezoidal rule
	double result = 0.0;
	for (std::size_t i = 0; i + 1 < N; ++i)
		result += 0.5 * (integrand[i] + integrand[i+1])
					  * (w_grid[i+1] - w_grid[i]);
	return result;
}

}  // namespace ikk
\end{lstlisting}

\noindent\rule{\linewidth}{0.4pt}

% ============================================================================
%  IMPORT SECTION J
%  Title: Open Variable Register
%  Status: Ongoing --- updated each draft pass
% ============================================================================
\section{Open Variable Register}
\label{sec:open-vars}

\noindent\rule{\linewidth}{0.8pt}
\textbf{Import block J} \hfill \textit{living document --- update each pass}
\noindent\rule{\linewidth}{0.4pt}

\begin{table}[h]
  \centering
  \begin{tabularx}{\linewidth}{l l X}
	\toprule
	Variable & Introduced & Status / constraint needed \\
	\midrule
	$\chi_w(w)$ & Eq.\ (F.5) &
	  Dark-scale metric correction. Formally undefined. Needs physical
	  derivation or phenomenological ansatz from clock-discrepancy data. \\
	$K_G(w)$ & Eq.\ (E.3) &
	  Gravitational projection kernel over dark coordinate. Known property:
	  $\int K_G\,dw = 1$. Explicit form unknown. \\
	$K_\text{EM}(w)$ & Eq.\ (E.3) &
	  EM projection kernel. Known property: truncates above $L_\text{EM}$.
	  Explicit form unknown. \\
	$\Pi_n$ & Eq.\ (B.3) &
	  Packing operator at scale $n$. Axiom only. Explicit construction
	  at each scale level needed. \\
	$\mathbf{T}_{AB}$ & Eq.\ (G.2) &
	  Trinary relation vector. Space defined. Dynamics not defined. \\
	$X_D$ & Table~\ref{tab:scale-ladder} &
	  Dark mediator. Placeholder only. Mass, spin, coupling unknown. \\
	$\hat{H}_w$ & Eq.\ (extended-TDSE) &
	  Dark-scale Hamiltonian. Without explicit form, the theory has
	  no new quantitative predictions at the QM level. \\
	$\mathbf{W}$ & Eq.\ (D.3) &
	  Hidden-intensity weight matrix. Must be PSD. Specific form
	  for each interaction channel needed. \\
	\bottomrule
  \end{tabularx}
  \caption{Open variables requiring further formalization.}
  \label{tab:open-vars}
\end{table}

\noindent\rule{\linewidth}{0.4pt}

\end{document}
